\section{Generative Paradigms for Novel View Synthesis}
\label{sec:generative}

The emergence of deep generative models---first GANs and then diffusion
models---has fundamentally altered the landscape of \nvs{}.
Where reconstruction-based methods infer a scene representation from observations
and render it deterministically, \emph{generative} methods learn a distribution
over plausible visual completions given partial observations.
This distinction is not merely architectural: it represents a shift from
\emph{interpolation} (rendering within the observed convex hull of viewpoints)
to \emph{extrapolation} (hallucinating plausible content beyond it).
Generative methods are therefore the primary vehicle for the one-image-to-3D,
text-to-3D, and single-image world-generation capabilities that bring \nvs{}
into the world-modeling regime.

\subsection{3-D-Aware GANs: Learning Geometry from 2-D Supervision}

The first wave of generative \nvs{} methods embedded radiance fields inside
GAN generators, enabling implicit 3-D supervision from unposed 2-D image
collections---a paradigm that proved influential precisely because it
required no 3-D ground truth.

GRAF~\citep{schwarz2020graf} pioneered this direction by conditioning a NeRF
generator on a latent shape-appearance code, training with a 2-D discriminator.
pi-GAN~\citep{chan2021pigan} improved the NeRF backbone with periodic activations
(SIREN~\citep{sitzmann2020siren}), achieving substantially better multi-view
consistency.
GIRAFFE~\citep{niemeyer2021giraffe} introduced scene compositionality by
combining foreground and background neural feature fields, enabling object-level
editing without explicit segmentation.
EG3D~\citep{chan2022efficient} is the landmark method of this era: its
\emph{triplane} representation---three axis-aligned 2-D feature maps decoded by
a shallow MLP---reduces rendering cost by orders of magnitude and admits a
super-resolution upsampler, enabling 512$\times$512 training at practical
throughput.
StyleNeRF~\citep{gu2022stylenerf} integrates the NeRF volume renderer into a
StyleGAN~\citep{karras2019stylegan2} backbone, inheriting its powerful
style-based latent structure for fine-grained appearance control.

A critical limitation of GAN-based approaches is \emph{mode coverage}: the
discriminator objective encourages sharp, realistic individual samples but does
not explicitly minimise distributional divergence over all plausible views,
making it prone to mode dropping for complex, asymmetric scenes.
This motivated the transition toward diffusion-based approaches.

\subsection{Diffusion Models for Camera-Controlled NVS}
\label{sec:diffusion_nvs}

Diffusion models~\citep{rombach2022ldm,ho2020ddpm} have become the
dominant generative paradigm for \nvs{} due to three decisive advantages:
(i)~training stability via a simple denoising objective;
(ii)~access to powerful image priors from large-scale pretraining;
(iii)~natural integration of arbitrary conditioning signals (pose, depth,
text, reference image) through classifier-free guidance.

\subsubsection{Pose-Conditioned Single-Image NVS}

The most direct application is \emph{single-image NVS}: given one reference
photograph and a target camera pose $\boldsymbol{\pi}^*$, generate a plausible
rendering at that viewpoint.
This is under-constrained without a prior---an infinite number of 3-D scenes
are consistent with a single 2-D observation---and diffusion models provide
exactly the distributional prior needed.

\paragraph{Zero-1-to-3 and successors.}
Zero-1-to-3~\citep{liu2023zero1to3} established the canonical formulation:
fine-tune a latent diffusion model (Stable Diffusion~\citep{rombach2022ldm})
on multi-view renders from Objaverse~\citep{deitke2023objaverse}, conditioning
on both the reference image and the relative pose $(\Delta\phi, \Delta\theta,
\Delta r)$ encoded as a CLIP-space vector.
The model learns to ``imagine'' how an object would look from any viewpoint,
without constructing an explicit 3-D representation.
Zero123++~\citep{shi2023zero123pp} strengthens multi-view consistency by
predicting a fixed grid of six canonical views simultaneously, reducing
inter-view inconsistency from independent sampling.
One-2-3-45~\citep{liu2023one2345} pipelines Zero-1-to-3 outputs into a
cost-volume-based MVS reconstruction~\citep{yao2018mvsnet}, converting the
diffusion output to a textured mesh within 45 seconds.

\paragraph{Synchronised multi-view denoising.}
A key challenge with independently conditioned diffusion is that each generated
view is sampled from a different realisation of the latent distribution,
producing inter-view geometric contradictions.
SyncDreamer~\citep{liu2023syncdreamer} addresses this by coupling the denoising
of multiple target views through shared 3-D volume attention at each diffusion
step, enforcing consistent geometry throughout the synthesis process.
MVDiffusion~\citep{tang2023mvdiffusion} constructs a correspondence-aware attention
mechanism that explicitly queries matching pixels across views, enabling
holistic panorama generation from a set of overlapping perspective images.
Wonder3D~\citep{long2024wonder3d} takes a two-domain approach: a cross-domain
diffusion model simultaneously denoises an RGB and a surface-normal image,
with the normal output serving as high-quality geometric supervision for
downstream mesh reconstruction.

\paragraph{Scalable and pose-free generation.}
EscherNet~\citep{kong2024eschernet} generalises pose conditioning via a
novel \emph{camera positional encoding}~(CaPE), enabling generation of
an arbitrary number of views at any resolution without per-angle fine-tuning.
ViVid-1-to-3~\citep{kwak2024vivid} recasts NVS as a video generation problem:
the temporal axis of a pre-trained video diffusion model is repurposed to model
the viewpoint axis, using the model's temporal consistency as a proxy for
multi-view coherence.
GeoNVS~\citep{chen2024geonvs} addresses geometric fidelity by conditioning a
video diffusion model on a rendered point cloud at the target pose, grounding
the generative hallucination in explicit 3-D geometry.

\subsubsection{Multi-View Diffusion and 3-D-Consistent Generation}

Single-view methods generate views independently or with limited cross-view
coupling.
A more ambitious goal is to train a diffusion model that \emph{natively}
produces multi-view consistent outputs---effectively encoding 3-D structure
implicitly in the model weights.

MVDream~\citep{shi2023mvdream} achieves this by replacing the temporal
self-attention in a video diffusion transformer with a 3-D-aware \emph{multi-view
attention} block that attends across all views simultaneously.
Training on rendered 3-D assets (Objaverse) with known camera poses teaches the
model to share structural information across views.
The resulting model can be used both as a standalone multi-view generator and
as a geometry-consistent score function for SDS-based text-to-3D.

GeNVS~\citep{chan2023genvs} takes a hybrid approach: a NeRF encoder constructs
a 3-D feature volume from the conditioning images, and the diffusion decoder
renders a target view by volumetrically querying this feature space.
This architecture provides hard geometric grounding while preserving the
expressiveness of the diffusion prior for high-frequency appearance details.

\subsubsection{Camera Trajectory Conditioning and Video Diffusion NVS}

A particularly important frontier is conditioning diffusion models on
\emph{continuous camera trajectories}, enabling the generation of smooth,
temporally coherent novel-view sequences---the direct bridge to world modeling.

\paragraph{Camera trajectory as a control signal.}
ViewCrafter~\citep{yu2024viewcrafter} conditions a video diffusion model
(CogVideoX) on point-cloud renderings along a target camera path, using
the rendered point cloud as a structural scaffold that prevents temporal
drift while the diffusion model fills in appearance and fine geometry.
Free3D~\citep{zheng2023free3d} injects camera pose into each frame's
cross-attention query, enabling a single video diffusion model to follow
arbitrary camera trajectories from a single reference image.
Novel View Extrapolation~\citep{shi2024nve} specifically targets
\emph{large-angle extrapolation}---the regime where geometric hallucination
is most necessary---using video diffusion priors to maintain perceptual
plausibility when the reference content provides little direct evidence.

\paragraph{Training-free and zero-shot camera control.}
NVS-Solver~\citep{you2024nvssolver} demonstrates that an off-the-shelf video
diffusion model already encodes sufficient geometric knowledge to solve
NVS in a zero-shot manner: by formulating view synthesis as a constrained
diffusion process with known-region guidance, it achieves competitive performance
without any pose-specific fine-tuning.
DDIM Inversion for NVS~\citep{xu2024ddim_nvs} exploits the deterministic
reconstruction property of DDIM~\citep{song2020ddim}: inverting the reference
image to its latent code and re-decoding with a modified pose conditioning
produces structurally faithful novel views, again without additional training.

\paragraph{Geometric conditioning for high fidelity.}
Projected Representation Conditioning~\citep{han2024prc} proposes projecting
3-D features from the reference view onto the target-view image plane,
providing an explicit geometric ``warp'' prior to the diffusion decoder.
This projected feature conditioning substantially reduces hallucination artifacts
in regions visible from the target but occluded in the reference.
A Single Image and Multimodality~\citep{tang2024smm} augments pose conditioning
with depth, surface normal, and semantic maps, each derived from monocular
estimation, to enrich the conditioning signal across multiple complementary
geometric cues.

\subsubsection{Reconstruction--Diffusion Hybrids}

A complementary design philosophy \emph{separates} the diffusion and reconstruction
stages: diffusion fills in missing views and hallucinations, while a deterministic
reconstruction pipeline (NeRF or 3DGS) enforces 3-D consistency.

Cat3D~\citep{gao2024cat3d} exemplifies this pattern: given sparse reference
views ($n\geq1$), a multi-view diffusion model generates a dense set of
synthetic pseudo-views, which are then fed to a standard 3DGS reconstruction
pipeline.
Because the diffusion model is fine-tuned jointly with the reconstruction loss,
the generated views are explicitly optimised for downstream reconstruction fidelity.
ReconFusion~\citep{wu2024reconfusion} integrates the diffusion prior more tightly:
it conditions a pixel-diffusion model on pixelNeRF~\citep{yu2021pixelnerf}
feature queries, acting as a learned regulariser that fills in uncertain regions
during sparse-view NeRF optimisation.
ZeroNVS~\citep{sargent2023zeronvs} handles the challenging case of
in-the-wild single images (non-object-centric, background present) by training
on scene-level data and conditioning on the camera intrinsics alongside pose,
enabling 360$^\circ$ generation from a single casual photograph.
NVComposer~\citep{li2024nvcomposer} addresses the case where multiple reference
views are available but unposed: a feature composition network aggregates
information from all reference images before passing to the diffusion model,
improving quality without requiring explicit pose estimation.

\subsection{Score Distillation: Diffusion as a 3-D Optimisation Oracle}

Score Distillation Sampling (SDS)~\citep{poole2023dreamfusion} opened a
separate pathway: instead of directly generating novel views, use a
pre-trained 2-D diffusion model as a differentiable \emph{score function}
to optimise a 3-D representation.
The core identity is:
\begin{equation}
  \nabla_\theta \mathcal{L}_\text{SDS}
  = \mathbb{E}_{t,\boldsymbol{\epsilon}}\!\left[
    w(t)\!\left(
      \hat{\boldsymbol{\epsilon}}_\phi(\mathbf{x}_t; y, t)
      - \boldsymbol{\epsilon}
    \right)\frac{\partial \mathbf{x}}{\partial \theta}
  \right],
  \label{eq:sds}
\end{equation}
where $\theta$ parameterises the 3-D representation, $\mathbf{x}$ is a rendered
image, $y$ is a text prompt, and $\hat{\boldsymbol{\epsilon}}_\phi$ is the
diffusion model's noise prediction.

DreamFusion~\citep{poole2023dreamfusion} first demonstrated that SDS can
synthesise coherent 3-D objects from text descriptions.
Magic3D~\citep{lin2023magic3d} introduces a coarse-to-fine strategy that
first optimises a coarse NeRF then refines a differentiable mesh, achieving
$8\times$ higher resolution.
ProlificDreamer~\citep{wang2023prolificdreamer} identifies the ``over-saturation''
artifact of SDS as a consequence of mode-seeking behaviour, proposing
\emph{Variational Score Distillation} (VSD) which models the 3-D parameter
distribution as a particle set and applies Stein variational gradient descent
for diverse, high-quality generation.
Fantasia3D~\citep{chen2023fantasia3d} disentangles geometry (optimised via
SDS on a DMTet mesh) from appearance (a physically-based BRDF), enabling
photorealistic materials and independent editing of shape and texture.
GaussianDreamer~\citep{yi2024gaussiandreamer} replaces NeRF with 3DGS as the
SDS target, leveraging Gaussian densification for fine detail and achieving
real-time rendering of the final asset.
LucidDreamer~\citep{chung2023luciddreamer} addresses SDS instability with an
interval-based score matching scheme that reduces gradient variance.

\subsection{Dedicated Multi-View Diffusion Architectures}

The methods above predominantly adapt existing video or image diffusion models.
A complementary line of work designs architectures \emph{from scratch} for
multi-view consistency.

MVDream~\citep{shi2023mvdream} replaces temporal attention with multi-view
attention and trains on Objaverse with ground-truth camera poses, producing
natively 3-D-consistent outputs.
ImageDream~\citep{wang2024imagedream} extends MVDream to image-conditioned
generation, enabling single-image-to-3D with strong semantic fidelity.
Zero123-XL~\citep{zero123xl2023} scales the Zero-1-to-3 paradigm to larger
model capacities and more diverse training data.
Direct3D~\citep{wu2024direct3d} trains a native 3-D latent diffusion model
on a triplane latent space, enabling direct 3-D generation without a two-stage
pipeline.

\paragraph{High-resolution and dense multi-view generation.}
ERA3D~\citep{li2024era3d} addresses the resolution bottleneck of earlier
multi-view diffusion models through a row-wise attention mechanism that
efficiently scales to higher-resolution output without the quadratic cost
of full 2-D attention.
MVDiffusion++~\citep{tang2024mvdiffusionpp} pushes further by using a dense
high-resolution attention design that couples overlapping views through
a learned correspondence graph, enabling faithful reconstruction from as
few as one or two input images.
SplatFlow~\citep{go2024splatflow} replaces the denoising diffusion process
with a multi-view rectified flow formulation, improving both training
efficiency and inference speed while maintaining strong 3-D consistency
for direct Gaussian splatting synthesis.

\paragraph{Large-scale 3D asset generation.}
Hunyuan3D~\citep{yang2024hunyuan3d} and its successor
Hunyuan3D~2.0~\citep{zhao2025hunyuan3dv2} combine a multi-view diffusion model
with a high-fidelity texture synthesis pipeline, scaling text-to-3D and
image-to-3D generation to production-quality textured meshes.
Bootstrap3D~\citep{sun2024bootstrap3d} addresses the data scarcity problem by
synthesising a large-scale multi-view image dataset with a self-training
loop, substantially improving the diversity and quality of multi-view
diffusion models trained on limited real data.
4Diffusion~\citep{zhang2024_4diffusion} extends the multi-view paradigm to
the temporal dimension, jointly modelling multiple views and multiple
timesteps in a single unified 4-D diffusion model.

\paragraph{Video diffusion as a 4-D generation backbone.}
Stable Video Diffusion~(SVD)~\citep{blattmann2023svd} provides the
pre-trained temporal prior that downstream 4-D methods build upon.
SV3D~\citep{voleti2024sv3d} fine-tunes SVD on orbital camera trajectories
to generate multi-view consistent 3-D objects from a single image.
SV4D~\citep{xie2024sv4d} extends this to dynamic objects by jointly conditioning
on camera trajectory and time, synthesising a consistent 4-D video from
a single reference.
AnimateDiff~\citep{guo2023animatediff} provides a plug-and-play temporal
module for personalised text-to-image models, enabling video generation
without per-model video training---a pattern widely adopted in subsequent
camera-controlled video diffusion work.
CogVideoX~\citep{yang2024cogvideox} scales video generation to high resolution
with an expert transformer architecture; its temporal attention blocks have
been adopted as a backbone by several NVS methods including
ViewCrafter~\citep{yu2024viewcrafter} and DimensionX~\citep{sun2024dimensionx}.
The Kling model family~\citep{kling2025report} and
Wan~\citep{wan2025wan} represent the latest generation of
commercial/open video foundation models that provide strong motion priors
for camera-controlled scene synthesis.

Table~\ref{tab:generative_methods} provides a structured comparison of generative
\nvs{} methods across the four paradigms above.

\begin{table}[t]
  \centering
  \caption{Generative NVS methods. \textbf{Backbone}: LDM=latent diffusion,
    VDM=video diffusion, GAN=generative adversarial.
    \textbf{Ctrl}: pose conditioning mechanism.
    \textbf{3D Out}: whether the method produces an explicit 3-D representation.}
  \label{tab:generative_methods}
  \small
  \setlength{\tabcolsep}{2pt}
  \begin{tabular}{@{}llllccc@{}}
    \toprule
    \textbf{Method} & \textbf{Venue}
      & \textbf{Backbone} & \textbf{Ctrl}
      & \textbf{3D Cons.} & \textbf{3D Out} & \textbf{Key Contribution} \\
    \midrule
    \multicolumn{7}{@{}l}{\textit{GAN-based}} \\
    GRAF~\citep{schwarz2020graf}                & NeurIPS'20 & NeRF-GAN   & Latent  & \pt & NeRF   & Radiance field GAN \\
    pi-GAN~\citep{chan2021pigan}                & CVPR'21    & NeRF-GAN   & Latent  & \pt & NeRF   & Periodic implicit \\
    GIRAFFE~\citep{niemeyer2021giraffe}         & CVPR'21    & Feat. Field& Latent  & \pt & Neural & Compositional \\
    EG3D~\citep{chan2022efficient}              & CVPR'22    & Triplane   & Latent  & \yes& Tri.   & Triplane+SR \\
    StyleNeRF~\citep{gu2022stylenerf}           & ICLR'22    & StyleGAN   & Latent  & \pt & NeRF   & Style-based NeRF \\
    \midrule
    \multicolumn{7}{@{}l}{\textit{Pose-conditioned single-image diffusion}} \\
    Zero-1-to-3~\citep{liu2023zero1to3}         & ICCV'23    & LDM        & $\Delta$pose & \pt & -- & Pose-cond LDM \\
    Zero123++~\citep{shi2023zero123pp}          & arXiv'23   & LDM        & $\Delta$pose & \pt & -- & 6-view grid \\
    One-2-3-45~\citep{liu2023one2345}           & NeurIPS'23 & LDM+MVS    & $\Delta$pose & \yes& Mesh & Diff$\to$MVS \\
    SyncDreamer~\citep{liu2023syncdreamer}      & ICLR'24    & LDM        & $\Delta$pose & \yes& -- & Volume sync \\
    Wonder3D~\citep{long2024wonder3d}           & CVPR'24    & LDM        & $\Delta$pose & \yes& -- & RGB+Normal \\
    EscherNet~\citep{kong2024eschernet}         & CVPR'24    & LDM        & CaPE    & \yes& -- & Arbitrary \#views \\
    ViVid-1-to-3~\citep{kwak2024vivid}          & arXiv'24   & VDM        & $\Delta$pose & \pt & -- & Video=NVS axis \\
    \midrule
    \multicolumn{7}{@{}l}{\textit{Camera trajectory / video diffusion}} \\
    Free3D~\citep{zheng2023free3d}              & ICLR'24    & VDM        & Cam embed& \pt & -- & Per-frame pose \\
    ViewCrafter~\citep{yu2024viewcrafter}       & arXiv'24   & VDM        & Ptcloud  & \yes& -- & Scaffold+VDM \\
    NVS-Solver~\citep{you2024nvssolver}         & arXiv'24   & VDM        & Zero-shot& \pt & -- & Guided diffusion \\
    DDIM Inversion~\citep{xu2024ddim_nvs}       & arXiv'24   & LDM        & Inversion& \pt & -- & Training-free \\
    GeoNVS~\citep{chen2024geonvs}               & arXiv'24   & VDM        & Ptcloud  & \yes& -- & Geo-grounded VDM \\
    Proj. Repr. Cond.~\citep{han2024prc}        & arXiv'24   & LDM        & 3D proj. & \yes& -- & Feature warp \\
    Novel View Extrap.~\citep{shi2024nve}       & arXiv'24   & VDM        & Cam traj.& \pt & -- & Large-angle extrap.\\
    \midrule
    \multicolumn{7}{@{}l}{\textit{Reconstruction--diffusion hybrids}} \\
    GeNVS~\citep{chan2023genvs}                 & ICLR'23    & LDM        & NeRF feat& \yes& NeRF & NeRF-cond diff. \\
    ZeroNVS~\citep{sargent2023zeronvs}          & CVPR'24    & LDM        & Pose+int.& \pt & -- & In-the-wild \\
    ReconFusion~\citep{wu2024reconfusion}       & CVPR'24    & LDM        & pixelNeRF& \yes& NeRF & Diff reg. \\
    Cat3D~\citep{gao2024cat3d}                  & NeurIPS'24 & MV-Diff    & $\Delta$pose & \yes& 3DGS& MV gen+recon \\
    NVComposer~\citep{li2024nvcomposer}         & arXiv'24   & LDM        & Feature  & \yes& -- & Unposed multi-ref\\
    \midrule
    \multicolumn{7}{@{}l}{\textit{Score distillation (text-to-3D)}} \\
    DreamFusion~\citep{poole2023dreamfusion}    & ICLR'23    & LDM+NeRF   & SDS      & \pt & NeRF & SDS \\
    Magic3D~\citep{lin2023magic3d}              & CVPR'23    & LDM+Mesh   & SDS      & \pt & Mesh & C2F SDS \\
    ProlificDreamer~\citep{wang2023prolificdreamer}& NeurIPS'23& LDM+NeRF & VSD      & \pt & NeRF & VSD \\
    Fantasia3D~\citep{chen2023fantasia3d}       & ICCV'23    & LDM+Mesh   & SDS      & \pt & BRDF & Disentangled \\
    GaussianDreamer~\citep{yi2024gaussiandreamer}& CVPR'24   & LDM+3DGS   & SDS      & \pt & 3DGS & SDS+3DGS \\
    Taming SDS~\citep{wang2024taming_sds}       & CVPR'24    & LDM+NeRF   & SDS      & \pt & NeRF & Mode collapse fix \\
    Text-to-3D CSD~\citep{yu2023text3dclassifier}& ICLR'24  & LDM+NeRF   & CSD      & \pt & NeRF & Classifier SDS \\
    \midrule
    \multicolumn{7}{@{}l}{\textit{Large-scale feed-forward reconstruction}} \\
    LRM~\citep{hong2023lrm}                     & ICLR'24    & Transformer& Single   & \yes& Tripl.& Triplane LRM \\
    TriplaneGaussian~\citep{zou2024triplanegaussian}& CVPR'24 & Transformer& Single  & \yes& 3DGS  & Triplane+GS \\
    GS-LRM~\citep{zhang2024gslrm}              & ECCV'24    & Transformer& Multi    & \yes& 3DGS  & Scalable LRM \\
    Instant3D~\citep{li2023instant3d}           & ICLR'24    & LRM+MVDiff & Text     & \yes& 3DGS  & Text$\to$3D fast \\
    CRM~\citep{wang2024crm}                     & arXiv'24   & Conv+LRM   & Single   & \yes& Mesh  & Conv reconstruct \\
    \midrule
    \multicolumn{7}{@{}l}{\textit{High-res multi-view / production 3D}} \\
    ERA3D~\citep{li2024era3d}                   & NeurIPS'24 & LDM        & $\Delta$pose & \yes& -- & Row-wise attn \\
    MVDiffusion++~\citep{tang2024mvdiffusionpp} & ECCV'24    & LDM        & Dense    & \yes& -- & Dense HR MV \\
    SplatFlow~\citep{go2024splatflow}           & arXiv'24   & RF-Diff    & MV       & \yes& 3DGS & Rectified flow \\
    Hunyuan3D~\citep{yang2024hunyuan3d}        & arXiv'24   & LDM        & Text/Img & \yes& Mesh  & Production 3D \\
    \midrule
    \multicolumn{7}{@{}l}{\textit{Video diffusion 4-D generation}} \\
    SV3D~\citep{voleti2024sv3d}                 & ECCV'24    & SVD        & Orbit    & \yes& 3DGS  & 3D from SVD \\
    SV4D~\citep{xie2024sv4d}                   & arXiv'24   & SVD        & Orbit+T  & \yes& 4DGS  & Dynamic 3D \\
    CAT4D~\citep{wu2024cat4d}                   & NeurIPS'24 & VDM        & View+T   & \yes& 4D    & 4D from MV \\
    4Diffusion~\citep{zhang2024_4diffusion}     & arXiv'24   & VDM        & MV+T     & \yes& --    & 4D video diff \\
    \bottomrule
  \end{tabular}
\end{table}

\paragraph{Implications for World Modeling.}
Diffusion-based \nvs{} represents a critical advance along World Modeling
Axis~2 (sparse/single-image generalisation, \textbf{D2}) and provides
the generative prior essential for completing unobserved scene regions
(Axis~1, geometric consistency \textbf{D1}).
Camera trajectory conditioning (\S\ref{sec:diffusion_nvs}) directly bridges
into Axis~3 (temporal dynamics), as a camera trajectory can be viewed as
a special case of a time-ordered action sequence.
The remaining gap is \emph{physical action conditioning} (\textbf{D4, D5}):
existing diffusion NVS methods condition on camera parameters but not on
object-level physical interactions, which remains an open problem
addressed in \S\ref{sec:open}.

% ─────────────────────────────────────────────────────────────────────────────
